\chapter{热力学系统的平衡态及状态方程}
\problem{1.1}{在一个大气压下，水的冰点温度是 $t_i$，汽点温度是 $t_s$。一支未经校准的温度计，在这样的压强下，放在冰水混合物中，指示温度为 $t'_i$；放在沸腾的水中指示温度为 $t'_s$。假定温度计的刻度在 $t'_i$ 和 $t'_s$ 之间是等分的，试确定中间指示温度 $t'_p$ 时，真实温度是多少？}
\answerarea
\problem{1.2}{将水银温度计浸在冰水中，水银柱长度为 $4.0\rm\,cm$；浸在沸水中，水银柱长度为 $24.0\rm\,cm$。
\begin{enumerate}
\item 温度计浸在某种沸腾的化学溶液中时，水银柱的长度为 $25.4\rm\,cm$，试确定溶液的温度。
\item 在室温 $22.0\,^\circ\mathrm C$ 时，水银柱的长度为多少？
\end{enumerate}}
\answerarea
\clearpage
\problem{1.3}{用 $L$ 表示液体温度计中液柱的长度，并定义温标 $t^*$ 与 $L$ 之间的关系为 $t^*=a\ln L+b$，式中 $a,b$ 为常量。规定冰点为 $t_i^*=0^\circ$、汽点为 $t_s^*=100^\circ$。设在冰点时液柱的长度为 $L_i=5.0\rm\,cm$，在汽点时液柱的长度为 $L_s=25.0\rm\,cm$，试确定 $t_1^*=0^\circ$ 到 $t_2^*=10^\circ$ 之间的液柱的长度差以及 $t_1^*=90^\circ$ 到 $t_2^*=100^\circ$ 之间的液柱长度差。}
\answerarea
\problem{1.4}{定义温标 $t^*$ 与测温属性 $X$ 之间的关系为 $t^*=\ln(kX)$，式中 $k$ 为常量。
\begin{enumerate}
\item 设 $X$ 为定体稀薄气体的压强，并假定水的三相点温度 $t^*=273.16^\circ$，试确定温标 $t^*$ 与理想气体温标之间的关系。
\item 在温标 $t^*$ 中冰点和汽点各多少度？
\item 在温标 $t^*$ 中是否存在 0 度？
\end{enumerate}}
\answerarea
\clearpage
\problem{1.5}{用定体气体温度计测得冰点的理想气体温度为 $273.15\rm\,K$，试确定温度计内的气体压强在冰点与其在水的三相点的比值的极限值。}
\answerarea
\problem{1.6}{定压气体温度计测温泡内的气体在水的三相点和汽点时体积比 $V_{\rm tr}/V_s$ 的极限值是 0.732038，试确定汽点在理想气体温标下的值。}
\answerarea
\clearpage
\problem{1.7}{利用定体气体温度计测量某物质的沸点。原来测温泡在水的三相点时，其中压强 $p_{\rm tr}=500\rm\,mmHg$；当测温泡浸入待测物质中时，测得的压强值为 $p=734\rm\,mmHg$。当从测温泡中抽出一些气体，使 $p_{\rm tr}$ 减为 $100\rm\,mmHg$ 时，测得 $p=146.68\rm\,mmHg$。再从测温泡中抽出一些气体，使 $p_{\rm tr}$ 减为 $40\rm\,mmHg$ 时，测得 $p=68.68\rm\,mmHg$。试确定该物质的沸点在理想气体温标下的温度。}
\answerarea
\problem{1.8}{记某定体气体温度计的测温泡中的压强在水的三相共存状态的温度为 $p_{\rm tr}$，在某确定温度 $T$ 下，测温泡中的气体的压强为 $p$。现对一状态的温度进行多次测量，以得到较准确的结果。测量数据如下（单位均为 $\rm kPa$）。\[\begin{array}{c|rrrrr}p_{\rm tr}&66.66&50.00&33.33&16.67&8.34\\p&102.35&76.76&51.19&25.59&12.81\end{array}\]试确定该状态在理想气体温标下的温度。}
\answerarea
\clearpage
\problem{1.9}{设计一温度计，其固定标准点为水的冰点和汽点，其间间隔记为 $100^\circ$。
\begin{enumerate}
\item 若记体积固定情况下上述两状态的压强的比值为 $\lim_{p_i\to0}p_s/p_i=\gamma_s$，则冰点在理想气体温标下的温度为 $\theta_i=100/(\gamma_s-1)$。试证明，利用定体气体温度计测量水的冰点的相对误差与测量汽点的相对误差之间有关系 $\left|\mathrm d\theta_i/\theta_i\right|=3.73\left|\mathrm d\gamma_s/\gamma_s\right|$。
\item 若一状态的压强与冰点压强的比值为 $\lim_{p_i\to0}p/p_i=\gamma$，则该状态在理想气体温标下的温度为 $\theta=\gamma\theta_i$。试证明，这样测量的相对误差与测量汽点的相对误差之间有关系 $\left|\mathrm d\theta/\theta\right|=\left|\mathrm d\gamma/\gamma\right|+3.73\left|\mathrm d\gamma_s/\gamma_s\right|$。
\end{enumerate}}
\answerarea
\problem{1.10}{仿上题，设计温度计时，仅选取一个固定标准点，该点的温度在理想气体温标下的数值为一普适常量。试说明：若记体积固定情况下一状态的压强与标准点的压强的比值为 $\lim_{p_{\rm tr}\to0}p/p_{\rm tr}=\gamma$，则该状态的温度测量的相对误差为 $\left|\mathrm d\theta/\theta\right|=\left|\mathrm d\gamma/\gamma\right|$。}
\answerarea
\clearpage
\problem{1.11}{设有 $1\rm\,mol$ 物质的量的理想气体系统，其热平衡态由 $(p,V)$ 两个状态参量标记，试给出状态 $(p_0,V_0)$ 附近的一个温度函数。}
\answerarea
\problem{1.12}{推广上题，设有一均匀物质系统，其在一定范围内的热平衡态可由参量 $(x,y)$ 确定，试证明在系统的某个状态 $(x_0,y_0)$ 附近可以定义温度函数 $T=f(x,y)$。}
\answerarea
\clearpage
\problem{1.13}{推广上题，设有一均匀物质系统，其在一定范围内的热平衡态可由参量 $(x,y,z,\ldots,w)$ 确定，试证明在系统的某个状态 $(x_0,y_0,z_0,\ldots,w_0)$ 附近可以定义温度函数 $T=f(x,y,z,\ldots,w)$。}
\answerarea
\problem{1.14}{通常人们称单位质量物质的体积为该物质的比容（或比体积）。试确定氮在压强为 $0.5\rm\,kg/cm^2$、温度为 $27\,^\circ\mathrm C$ 情况下的比容。}
\answerarea
\clearpage
\problem{1.15}{一房间的容积为 $10\times10\times3\rm\,m^3$，试确定：在标准状况下，这间房子内的空气的质量为多少千克？已知空气的平均摩尔质量是 $29\times10^{-3}\rm\,kg/mol$。}
\answerarea
\problem{1.16}{一立方体形状的容器的内边长为 $20\rm\,cm$，其中贮有 $1.0\rm\,atm$、$300\rm\,K$ 的气体，当把气体加热到 $400\rm\,K$ 时，容器每个壁所受的压力为多大？}
\answerarea
\clearpage
\problem{1.17}{如图所示，两个截面积相同的连通管，一为开管，一为闭管，原来两管内的水银面等高，且闭管内水银面到管顶距离为 $a$。今打开活塞，使水银漏掉一些，因此开管内水银下降了 $h$，试确定闭管内水银面下降了多少？设大气压为 $p_0$，整个过程中温度保持不变。
\fig{1_17}}
\answerarea
\problem{1.18}{截面积为 $S$ 的粗细均匀的 U 形管，其中贮有水银，高度如图所示。今将 U 形管的右侧与大气相通，左侧上端封闭，其中空气柱的温度为 $300\rm\,K$。若要使封闭的管中的空气柱的长度变为 $60\rm\,cm$，需加热到多少 K？已知大气压强保持为 $75\rm\,cmHg$。
\fig{1_18}}
\answerarea
\clearpage
\problem{1.19}{一端封闭的玻璃管长 $l=70.0\rm\,cm$，贮有空气，气柱上面有一段长为 $h=20.0\rm\,cm$ 的水银柱将气柱封住。水银面与管口对齐。今将玻璃管的开口端用玻璃片盖住，轻轻倒转后再除去玻璃片，因而使一部分水银漏出。
\begin{enumerate}
\item 当大气压 $p_0=75.0\rm\,cmHg$ 时，留在管内的水银柱有多长？
\item 在什么样的大气压强下水银完全从管内流出？
\end{enumerate}}
\answerarea
\problem{1.20}{水银气压计内混进了一个空气泡，因此它的读数比实际的气压小。当精确的气压计读数为 $768\rm\,mmHg$ 时，它的读数只有 $748\rm\,mmHg$。此时管内水银面到管顶的距离为 $80\rm\,mm$，问当此气压计的读数为 $734\rm\,mmHg$ 时，实际大气压应是多少？假设空气的温度保持不变。}
\answerarea
\clearpage
\problem{1.21}{一个氢气球可以膨胀，随着气球不断升高，大气压强不断减小，氢气球就不断膨胀。如果忽略掉大气温度和摩尔质量随高度的变化，试问气球在上升过程中所受的浮力是否变化？}
\answerarea
\problem{1.22}{在矿井入风巷道的某一截面处空气的压强 $p=0.9\rm\,bar$，温度 $t=17\,^\circ\mathrm C$，流速 $v=5\rm\,m/s$，该处截面积 $S=8\rm\,m^2$，试确定每秒流经该处的空气有多少千克？（已知空气的平均摩尔质量是 $28.9\times10^{-3}\rm\,kg/mol$）}
\answerarea
\clearpage
\problem{1.23}{一个氧气瓶容积是 $32\rm\,L$，其中氧气的压强是 $130\rm\,atm$。规定瓶内氧气压强降到 $10\rm\,atm$ 时就得充气，以免混入其它气体而需洗瓶。今有一玻璃室，每天需用 $1.0\rm\,atm$ 的氧气 $400\rm\,L$，试确定一瓶氧气能用几天？}
\answerarea
\problem{1.24}{一篮球在室温为 $20\,^\circ\mathrm C$ 时打入空气，使其内压强达到 $1.5\rm\,atm$，试确定：
\begin{enumerate}
\item 赛球时，篮球温度升高到 $40\,^\circ\mathrm C$，这时球内的压强有多大？
\item 在球赛过程中，球被扎破了一个小洞，开始漏气。球赛结束后，篮球恢复到室温时，球内剩下的空气是原有空气的百分之几？（篮球体积不变，室内外气压均为 $1\rm\,atm$）
\end{enumerate}}
\answerarea
\clearpage
\problem{1.25}{在密封的瓶子中贮有温度为 $17\,^\circ\mathrm C$、压强为 $1\rm\,atm$ 的空气。如果瓶内压强是 $1.3\rm\,atm$ 时就可将瓶塞顶开，试确定，当瓶内空气被加热到多高温度时就密封不住了？}
\answerarea
\problem{1.26}{28 英寸自行车车轮直径为 $71.12\rm\,cm$，内胎截面半径为 $3.0\rm\,cm$，在 $-3\,^\circ\mathrm C$ 的天气里向空胎内打气。打气筒长为 $30\rm\,cm$，截面半径为 $2.5\rm\,cm$，打了 20 下之后，如果车胎内气体的温度为 $7\,^\circ\mathrm C$，试估算车胎内气体的压强。}
\answerarea
\clearpage
\problem{1.27}{深海潜水员要在四周都是水的压力下呼吸空气。因为在 $0.2\rm\,MPa$ 分压强下的氧是有毒的，所以在一定的水深以下必须使用特殊的气体混合物。已知海水的密度为 $\rho=1.025\times10^3\rm\,kg/m^3$，试确定：
\begin{enumerate}
\item 按照含氧 $21\%$ 的体积百分比计算，在什么深度下空气中氧的分压强等于 $0.2\rm\,MPa$？
\item 在深水作业中使用含 $3\%$ 的氧和 $97\%$ 的氦（体积百分比）的气体混合物，在水深 $200\rm\,m$ 时，这种气体混合物中氧的分压强是多少？
\end{enumerate}}
\answerarea
\problem{1.28}{按质量计，空气由 $76\%$ 的氮、$23\%$ 的氧和约 $1\%$ 的氩组成（其余组分很少，可以忽略），试确定空气的平均摩尔质量。}
\answerarea
\clearpage
\problem{1.29}{两容器容积相同，装有相同质量的氮气和氧气，用一内壁光滑的水平玻璃管相通，玻璃管正中间有一小滴水银，如图所示，试确定：
\begin{enumerate}
\item 如果两容器内气体的温度相同，水银滴能否平衡？
\item 如果将氧的温度保持为 $T_1=30\,^\circ\mathrm C$，氮的温度保持为 $T_2=0\,^\circ\mathrm C$，水银滴如何移动？
\item 要使水银滴不动，并维持两边温度差为 $30\,^\circ\mathrm C$，则氮的温度应为多高？
\end{enumerate}
\fig{1_29}}
\answerarea
\problem{1.30}{容积为 $2.5\rm\,L$ 的烧瓶内有 $1.0\times10^{15}$ 个氧分子、$4.0\times10^{15}$ 个氮分子和 $3.3\times10^{13}$ 个氩分子，设混合气体的温度为 $150\,^\circ\mathrm C$，试确定混合气体的压强。}
\answerarea
\clearpage
\problem{1.31}{试确定压强为 $100\rm\,atm$、密度为 $100\rm\,g/L$ 的氧气的温度，已知氧气的范德瓦尔斯修正量 $a=1.360\rm\,atm\cdot(L/mol)^2$，$b=0.031\,831\rm\,L/mol$。}
\answerarea
\problem{1.32}{试用范德瓦尔斯方程确定密闭于容器内质量 $M=1.1\rm\,kg$ 的 $\mathrm{CO_2}$ 的压强。已知容器的容积为 $V=20\rm\,L$，气体的温度 $t=13\,^\circ\mathrm C$。并将所得结果与用理想气体模型下的结果比较。（已知 $\mathrm{CO_2}$ 的范德瓦尔斯修正量 $a=3.592\rm\,atm\cdot(L/mol)^2$，$b=0.042\,671\rm\,L/mol$）}
\answerarea
\clearpage
\problem{1.33}{已知一热力学系统的状态方程为 $f(p,V,T)=0$，试证明，对于一般的函数都有 $\displaystyle\left(\frac{\partial V}{\partial p}\right)_T\left(\frac{\partial p}{\partial T}\right)_V\left(\frac{\partial T}{\partial V}\right)_p=-1$。}
\answerarea
\problem{1.34}{试分别对理想气体和范德瓦尔斯气体，确定其体膨胀系数、等体压强系数和等温压缩系数。}
\answerarea
\clearpage
\problem{1.35}{实验测得，常温下水银的体膨胀系数为 $1.80\times10^{-4}\rm/K$，等温压缩系数为 $3.87\times10^{-6}\rm/atm$。并且，在温度为 $278\rm\,K$、压强为 $2.0\rm\,atm$ 情况下，水银的密度为 $13.60\times10^3\rm\,kg/m^3$。对水银加压并冷却，在温度为 $268\rm\,K$ 情况下测得其密度为 $13.75\times10^3\rm\,kg/m^3$，试确定该情况下水银的压强。}
\answerarea
\problem{1.36}{实验测得，某种液体在常温下的等温压缩系数近似为常量 $4.5\times10^{-5}\rm/atm$，并且，在 $0.0\,^\circ\mathrm C$ 温度下的压强为 $1.0\rm\,atm$，在 $100.0\,^\circ\mathrm C$ 温度下的压强为 $2.0\rm\,atm$。试确定该液体在常温下的体膨胀系数。}
\answerarea
\clearpage
\problem{1.37}{对一维系统，仿照体膨胀系数，我们可以定义线膨胀系数。对厚度都为 $d$、长度和宽度也都分别相同的两个金属薄片测得其线膨胀系数分别为 $\alpha_1,\alpha_2$，为方便标记，记 $\alpha_1<\alpha_2$。现将它们粘成厚度为 $2d$ 的薄板，当温度升高 $\Delta T$ 时，该薄板将变成球壳状。试确定该球壳的曲率半径。}
\answerarea
\problem{1.38}{测量标定，一以钢丝为摆线的摆钟在 $18\,^\circ\mathrm C$ 条件下是准确的。如果该钢丝的线膨胀系数在常温下为常量 $1.2\times10^{-5}/^\circ\mathrm C$，那么在温度为 $28\,^\circ\mathrm C$ 的情况下该摆钟每天差多少秒？}
\answerarea
\clearpage
\problem{1.39}{实验测量表明，固体和液体通常不易被压缩，并且它们的体膨胀系数 $\alpha$ 和等温压缩系数 $\kappa_T$ 都是常量，试确定它们的状态方程。}
\answerarea
\problem{1.40}{实验上测得某种气体在压强为 $p$、体积为 $V$、温度为 $T$ 的状态附近的等体压强系数为 $\beta=1/T$，等温压缩系数为 $\kappa_T=\dfrac1p\left(1-\dfrac bV\right)$，其中 $b$ 为数值大于零的常量，试确定该气体系统的状态方程。}
\answerarea
\clearpage
\problem{1.41}{实验测得一 $p$-$V$-$T$ 系统的体膨胀系数和等温压缩系数与状态参量的关系分别为 $\alpha=3aT^3/V$，$\kappa_T=b/V$，其中 $a$ 和 $b$ 为常量，试确定该系统的状态方程。}
\answerarea
\problem{1.42}{实验测得 $1\rm\,mol$ 过热蒸气的体膨胀系数和等温压缩系数与状态参量的关系分别为 $\alpha=\dfrac{R}{pV_m}+\dfrac{na}{V_mT^{n+1}}$，$\kappa_T=\dfrac{RT}{p^2V_m}$，其中 $n,a$ 和 $R$ 为常量，试确定该过热蒸气系统的状态方程（气体状态方程的卡兰达模型）。}
\answerarea
\clearpage
\problem{1.43}{近代物理学中常用电子伏（eV）作为能量单位。试确定在多高温度下分子的平均平动动能为 $1\rm\,eV$？$1\rm\,K$ 温度的热运动平均能量相当于多少电子伏？}
\answerarea
\problem{1.44}{对一致密天体测得其温度为 $3.87\times10^8\,^\circ\mathrm C$。假设该星体由核子（N）、$\Delta$ 共振态粒子及 $\Lambda$ 超子组成，试确定这些粒子的热运动动能各为多少焦耳。实验测得这些粒子的质量分别是 $m_N=938\rm\,MeV$，$m_\Delta=1232\rm\,MeV$，$m_\Lambda=1115\rm\,MeV$，试估计不考虑温度效应情况下对该星体物质的性质进行计算带来的系统误差为多大。}
\answerarea
\clearpage
\problem{1.45}{一能量为 $1.0\rm\,TeV$ 的高能宇宙射线粒子射入 $0.001\rm\,mol$ 的氖管中。如果该宇宙射线粒子的能量完全由氖气吸收，试确定平衡后氖气的温度升高多少？}
\answerarea
\problem{1.46}{我们知道，聚变能可以抽象为氘原子核与氘原子核融合形成氦原子核时释放出的能量，但是，由于氘核带正电荷，其间有库仑排斥力，因此需要温度高到一定程度才可能有聚变反应发生。实验测得氘原子核的有效直径大约为 $3\rm\,fm$，试确定该聚变反应的点火温度。}
\answerarea
\clearpage
\problem{1.47}{一容积为 $11.2\rm\,L$ 的真空系统已被抽到 $1.0\times10^{-5}\rm\,mmHg$ 的真空，为了提高其真空度，将它放入 $300\,^\circ\mathrm C$ 的烘箱烘烤，使器壁上吸附的气体分子释放出来。若烘烤后压强增为 $1.0\times10^{-2}\rm\,mmHg$，问器壁上原来吸附了大约多少个气体分子？}
\answerarea
\problem{1.48}{质量为 $50.0\rm\,g$、温度为 $18.0\,^\circ\mathrm C$ 的氦气装在容积为 $10.0\rm\,L$ 的封闭容器内，容器以 $v=200\rm\,m/s$ 的速度作匀速直线运动。若容器突然停止，定向运动的动能全部转化为分子热运动的动能，则平衡后氦气的温度和压强将各增大多少？}
\answerarea
\clearpage

